\section{索引与复现}

\begin{verbatim}
npm start                 # 启动，默认 127.0.0.1:8765
npm test                  # 257 项自动测试
npm run test:smoke        # 无头浏览器真打一局
npm run build:knowledge   # 从引擎重新生成知识卡
npm run eval:retrieval    # 检索四臂对照
\end{verbatim}

正文数字的来源：规模与测试数取自 \code{npm test} 的现场输出；检索对照、工具决策三轮、平衡矩阵与对手 agent 实测分别见 \path{reports/} 下的 \code{retrieval-extended.json}、\code{live-model-eval*.json}、\code{balance-matrix.json} 与 \code{opponent-agent.json}。宠物数、属性数与知识卡数都能用一行命令现场复算。

\textbf{交付与发布。} 仓库目前是本机工程。我会在提交后整理成一个可独立运行的仓库发布到 GitHub：去掉与本机环境绑定的脚本路径，补齐 README 的启动步骤与依赖说明，使克隆之后执行 \code{npm start} 即可完整游玩、\code{npm test} 即可复现全部结论；只有模型实验需要自己的密钥。
